\begin{dedicatoria}
   \vspace*{\fill}
   \centering
   \noindent
   \textit{ À criança curiosa que um dia me habitou e que, apesar de todos os anos e de todas as respostas, nunca deixou de perguntar: "por quê?". É essa mesma curiosidade, hoje lapidada pelo método e pelo rigor científico, que amadureceu o olhar de quem um dia observava o mundo com espanto e hoje busca compreendê-lo por meio da ciência.} \vspace*{\fill}
\end{dedicatoria}

% --------------------------------------------------------------------------------------------
%										Agradecimentos
% --------------------------------------------------------------------------------------------

\begin{agradecimentos}

Quero expressar minha profunda gratidão a todos que contribuíram, de alguma forma, para esta jornada de aprendizado, crescimento e superação ao longo do mestrado.

Agradeço, primeiramente, à minha família, em especial aos meus pais, Evandro Costa Cavalcante e Ivaneide Silva de Brito Cavalcante, pelo amor incondicional, apoio constante e incentivo em todos os momentos da minha vida. Agradeço também à minha irmã, Jessica Costa Cavalcante, e ao meu sobrinho, Gutierréz Cavalcante, por serem uma fonte constante de alegria e inspiração.

Ao meu amigo de mestrado, Jonas Leite, pelo companheirismo, pelas discussões enriquecedoras, pelo apoio nos momentos de dificuldade e pelas risadas compartilhadas ao longo dessa jornada.

Agradeço aos meus amigos Frank Marcelo, Naiane Barreto, Márcia Barbosa, Ligia Cristina, Jailany Araujo, Joaria Silva, Karla Pereira, Maria Suêd, Eduarda Monteiro, Erica Monteiro e Elayne Monteiro, pela amizade, apoio e incentivo durante todo o percurso do mestrado. Cada um de vocês contribuiu de maneira única para tornar essa experiência mais significativa e memorável.

Agradeço ao meu orientador, Prof. Dr. Luciano Barosi de Lemos, pela orientação, paciência e incentivo durante todo o desenvolvimento desta dissertação. Agradeço também aos demais professores do Programa de Pós-Graduação em Física da Universidade Federal de Campina Grande, que contribuíram para minha formação acadêmica e científica.

Por fim, agradeço à Coordenação de Aperfeiçoamento de Pessoal de Nível Superior (CAPES) – Brasil – Código de Financiamento 001, pelo apoio financeiro que possibilitou a realização deste trabalho.
\end{agradecimentos}

% --------------------------------------------------------------------------------------------
%										Epígrafe
% --------------------------------------------------------------------------------------------

\begin{epigrafe}
	\vspace*{\fill}
		\begin{flushright}
			\textit{"A natureza só utiliza os fios mais longos para tecer seus padrões, \\
			        mas cada pequena parte de seu tecido revela \\
					a organização de toda a tapeçaria." (Richard P. Feynman)}
		\end{flushright}
\end{epigrafe}